% Paper 2 of 7 in The Windstorm Series
\documentclass[12pt,a4paper]{article}

\usepackage{amsmath}
\usepackage{amssymb}
\usepackage{graphicx}
\usepackage{hyperref}
\usepackage{natbib}
\usepackage{booktabs}

\title{Rate-Distortion Bounds on Serial Decoding Throughput}

\author{Grant Lavell Whitmer III\\
The Windstorm Institute, Fort Ann, NY 12827, USA\\
\texttt{grantwhitmer3@gmail.com}}

\date{March 2026}

\begin{document}

\maketitle

\begin{abstract}
This study derives a mechanistic bound on effective information per serial decoding event using the $M$-ary rate-distortion function $R_M(\varepsilon)$. Applied to three independently characterized systems---the ribosome ($M = 21$, $\varepsilon = 10^{-4}$), human phonology ($M = 31$, $\varepsilon \sim 0.02$), and the chromatic scale ($M = 12$, $\varepsilon \sim 0.05$)---the framework predicts throughput values of 4.39, 4.71, and 3.12 bits, respectively. Under realistic non-uniform amino acid frequencies, the ribosome prediction shifts from 4.390 to approximately 4.175 bits, revealing a 0.2-bit excess-capacity margin. The vocabulary-independence prediction was tested across 1,749 models spanning two architectures---32 causal language models and 1,717 encoder-decoder translation models---evaluated on a shared reference corpus. Vocabulary size does not predict bits-per-byte in causal models ($\beta_1 = 0.058$, $p = 0.643$) or in language-matched translation models (Spearman $p = 0.13$). Language-mismatch stratification reveals a quantifiable $+1.33$-bit information cost when models process unfamiliar languages. These results support the hypothesis that the rate-distortion floor is a universal mechanistic bound for serial decoders and vocabulary size is primarily a redundancy parameter.
\end{abstract}

\noindent\textbf{Keywords:} $M$-ary discrimination; receiver constraint; genetic code fidelity; tokenizer vocabulary; bits per byte; cross-architecture compression; language mismatch penalty; information bottleneck; throughput convergence; substrate independence

\section{Introduction}

In a companion paper~\cite{whitmer2026fons}, we predicted that AI tokenizer vocabularies would cluster near 64 symbols, mirroring the genetic code. That prediction was falsified: AI tokenizers use 30K--256K tokens. However, the falsification revealed that while vocabulary sizes differ by orders of magnitude, the effective information per sequential processing event occupies a comparatively narrow range across substrates.

\subsection{Scope and Definitions}

This paper concerns serial discrimination channels---systems where a receiver decodes one symbol per time step under noise. Three quantities are distinguished:

\textbf{Definition 1 (Support-size upper bound).}
\begin{equation}
    H_0(X) = \log_2 |X|
\end{equation}

\textbf{Definition 2 (Source entropy rate).}
\begin{equation}
    H(X_t \mid C_t)
\end{equation}

\textbf{Definition 3 (Receiver-resolved information).}
\begin{equation}
    I_{\mathrm{eff}} = I(X_t; Y_t \mid C_t)
\end{equation}

These obey: $I(X_t; Y_t \mid C_t) \leq H(X_t \mid C_t) \leq \log_2 |X|$. The constraint does not apply to parallel visual processing, continuous analog signals, or systems without serial discrimination.

\section{Materials and Methods}

\subsection{Rate-Distortion Floor Derivation}

For a source with $M$ equiprobable symbols and symmetric error rate $\varepsilon$, the $M$-ary rate-distortion function gives the minimum information per event for reliable reconstruction:
\begin{equation}
    R_M(\varepsilon) = \log_2(M) - H_b(\varepsilon) - \varepsilon \log_2(M - 1)
\end{equation}
where $H_b(\varepsilon) = -\varepsilon \log_2(\varepsilon) - (1 - \varepsilon) \log_2(1 - \varepsilon)$ is the binary entropy function~\cite{shannon1948}.

\subsection{Application to Three Biological Systems}

The rate-distortion floor was applied to three independently characterized serial decoding systems:

\begin{enumerate}
    \item \textbf{Ribosome} ($M = 21$ amino acids, $\varepsilon = 10^{-4}$): The error rate reflects the combined fidelity of aminoacyl-tRNA synthetase selection and ribosomal proofreading~\cite{zaher2009}.
    \item \textbf{Human phonology} ($M = 31$ cross-linguistic median phonemes, $\varepsilon \sim 0.02$): Error rate estimated from confusion matrices in noisy listening conditions~\cite{millernicely1955}.
    \item \textbf{Chromatic scale} ($M = 12$ tones, $\varepsilon \sim 0.05$): Error rate estimated from pitch discrimination thresholds in musical perception.
\end{enumerate}

\subsection{Robustness Analysis: Non-Uniform Symbol Distributions}

For a source with non-uniform distribution $p_i$ over $M$ symbols, the mutual information under a symmetric error channel becomes:
\begin{equation}
    I(X;Y) = H(\mathbf{p}) - H_b(\varepsilon) - \varepsilon \log_2(M - 1) + \varepsilon \, D_{\mathrm{KL}}(\mathbf{p} \| \mathrm{uniform}) \cdot O(\varepsilon)
\end{equation}
where $H(\mathbf{p}) = \sum p_i \log_2(1/p_i)$ replaces $\log_2(M)$. For the ribosome, empirical amino acid frequencies from UniProt yield $H(\mathbf{p}) \sim 4.177$ bits versus $H(\mathrm{uniform}) = 4.392$ bits---a 0.215-bit reduction~\cite{zaher2009}.

\subsection{Tokenizer Sweep: 1,749 Models}

\subsubsection{Model Selection}

Two architecture classes were evaluated:
\begin{itemize}
    \item \textbf{32 causal language models} (GPT-2, LLaMA, Pythia, Mistral, and others) spanning vocabulary sizes of 30K--250K tokens.
    \item \textbf{1,717 encoder-decoder translation models} (Helsinki-NLP/OpusMT) spanning vocabulary sizes of 3,932--80,379 tokens.
\end{itemize}

\subsubsection{Evaluation Metric}

Bits-per-byte (BPB) normalizes across tokenizations and architectures. All models were evaluated using teacher-forced cross-entropy on WikiText-2 (English)~\cite{berger1971}.

\subsubsection{Language-Match Stratification}

Because translation models were trained on specific language pairs but evaluated on English, they were separated into three groups by language match:
\begin{itemize}
    \item \textbf{EN$\to$X (matched encoder, $n = 157$):} English input matches corpus; decoder outputs foreign language.
    \item \textbf{X$\to$EN (matched decoder, $n = 256$):} Foreign input mismatches corpus; decoder trained for English output.
    \item \textbf{X$\to$Y (fully mismatched, $n = 1{,}304$):} Neither encoder nor decoder expects English.
\end{itemize}

\subsubsection{Hardware}

All evaluations were performed on an NVIDIA RTX 5090 GPU with deterministic inference settings. The complete experiment protocol is provided in the supplementary materials~\cite{whitmer2026rlf_code}.

\subsubsection{Statistical Methods}

Vocabulary--BPB relationships were assessed using ordinary least squares regression ($\beta_1$ coefficient) for causal language models and Spearman rank correlation ($\rho$) for translation model groups. The coefficient of determination ($R^2$) quantified explained variance. Within-family analysis used the Pythia model series (14M to 1.4B parameters, identical tokenizer) to isolate the effect of model capacity from vocabulary size. Between-group differences were tested using the Kruskal--Wallis $H$ test.

\subsection{Independent Confirmation}

Tlusty~\cite{tlusty2008} derived $M \sim 20$--25 for the genetic code from the chromatic number theorem (Ringel--Youngs)---a completely independent topological argument (graph coloring versus rate-distortion) converging on the same biological parameter.

\section{Results}

\subsection{Rate-Distortion Predictions}

Table~\ref{tab:rd_predictions} presents the zero-free-parameter predictions from the rate-distortion floor applied to three biological systems.

\begin{table}[ht]
\centering
\caption{Rate-distortion predictions for three serial decoding systems.}
\label{tab:rd_predictions}
\begin{tabular}{lccccc}
\toprule
System & $M$ & $\varepsilon$ & $R_M(\varepsilon)$ & Observed & Residual \\
\midrule
Ribosome & 21 & $10^{-4}$ & 4.390 bits & 4.390 bits & $+0.000$ \\
Phonology & 31 & 0.02 & 4.71 bits & 4.95 bits & $+0.24$ \\
Chromatic scale & 12 & 0.05 & 3.12 bits & 3.58 bits & $+0.46$ \\
\bottomrule
\end{tabular}
\end{table}

At $\varepsilon = 10^{-4}$, the ribosome correction terms amount to only 0.0019 bits (0.04\% of $\log_2(21)$). The near-exact match between $R_M(\varepsilon)$ and $\log_2(M)$ at very low error rates is a mathematical property of Equation~(4), not an independent empirical confirmation. The framework's value lies in the non-trivial cases (phonology and chromatic scale) and in unifying $M$, $\varepsilon$, and throughput into a single testable relationship.

The residuals represent excess-capacity margins above the rate-distortion floor. Biology operates nearest the floor; cognition and music carry higher slack, consistent with greater discrimination cost in neural substrates.

\subsection{Non-Uniform Distribution Robustness}

Under non-uniform amino acid frequencies, the ribosome prediction shifts from 4.390 to approximately 4.175 bits. The observed ribosome throughput therefore sits approximately 0.2 bits above the non-uniform floor---a small but real excess-capacity margin consistent with the slack observed in other substrates. This correction makes the result more interesting rather than less: rather than a suspiciously perfect residual of 0.000, a genuine physical margin comparable to what Hopfield kinetic proofreading predicts is revealed~\cite{hopfield1974}.

\subsection{Vocabulary Independence Across 1,749 Models}

Table~\ref{tab:vocab_independence} presents the vocabulary-independence results across all four model groups.

\begin{table}[ht]
\centering
\caption{Vocabulary independence across four model groups.}
\label{tab:vocab_independence}
\begin{tabular}{lccccccl}
\toprule
Group & $n$ & BPB mean & BPB median & Statistic & $p$-value & $R^2$ & Vocab-indep.? \\
\midrule
Causal LMs & 32 & 0.94 & 0.85 & $\beta_1 = 0.058$ & 0.643 & 0.007 & Yes \\
EN$\to$X & 157 & 4.19 & 3.75 & $\rho = -0.122$ & 0.129 & 0.002 & Yes \\
X$\to$EN & 256 & 5.98 & 5.56 & $\rho = -0.068$ & 0.280 & 0.016 & Yes \\
X$\to$Y & 1,304 & 5.52 & 5.49 & $\rho = -0.068$ & 0.015 & 0.032 & Weak* \\
\bottomrule
\end{tabular}

\smallskip
\noindent\footnotesize{*Statistically significant at $p < 0.05$ but $\rho = -0.068$ indicates negligible practical effect; vocabulary explains only 3.2\% of variance.}
\end{table}

Vocabulary size does not predict BPB in causal language models ($p = 0.643$, $R^2 = 0.007$) or in language-matched translation models ($p = 0.129$, $R^2 = 0.002$). The vocabulary-independence result holds across both architectures and all four language-match conditions. Figures 1--4 illustrate these relationships (submitted separately).

\subsection{Within-Family Analysis}

Within the Pythia model family, which uses an identical tokenizer across all model sizes, BPB drops 46\% from 14M to 1.4B parameters. Model capacity drives compression; vocabulary does not.

\subsection{Language-Mismatch Finding}

The separation into language-match groups revealed an unexpected result: the information cost of processing mismatched input is quantifiable.

\begin{table}[ht]
\centering
\caption{Language-match effect on throughput.}
\label{tab:lang_mismatch}
\begin{tabular}{lccp{6cm}}
\toprule
Condition & $n$ & BPB mean & Interpretation \\
\midrule
EN$\to$X (matched encoder) & 157 & 4.19 & Model understands input; efficient processing \\
X$\to$Y (fully mismatched) & 1,304 & 5.52 & Model confused on both sides; $+1.33$ bits penalty \\
X$\to$EN (matched decoder) & 256 & 5.98 & Encoder garbled, decoder compensates; highest cost \\
\bottomrule
\end{tabular}
\end{table}

Matched-encoder models (EN$\to$X), which receive English input matching the corpus, operate at BPB $\sim 4.19$---within the throughput basin and near the ribosome prediction of 4.39 bits. Fully mismatched models (X$\to$Y) operate at BPB $\sim 5.52$---a penalty of $+1.33$ bits. The groups are significantly different (Kruskal--Wallis $H = 230.2$, $p \sim 0$), but vocabulary independence holds within each group. This finding was not part of the original experiment design and emerged during post-hoc analysis. It is reported as an exploratory finding, though the effect size (1.33 bits) is large and consistent across 1,461 models.

\subsection{Empirical Phonology Values}

Two phonology measurements appear in this series:
\begin{itemize}
    \item \textbf{Support-size bound [SS]:} $\log_2(31) = 4.95$ bits---the maximum entropy of the phoneme inventory assuming uniform distribution.
    \item \textbf{Confusion-matrix MI [MI]:} Miller and Nicely~\cite{millernicely1955} measured $I(X;Y) \sim 4.20$ bits for English consonants---a direct Tier~1 estimate of receiver-resolved information under noise.
\end{itemize}

The rate-distortion prediction $R_{31}(0.02) = 4.71$ bits sits between these two values. Measured against the support-size bound, the residual is $+0.24$ bits (below the ceiling, as expected). Measured against the confusion-matrix MI, the residual is $-0.51$ bits (the floor overpredicts, suggesting that contextual redundancy in real speech reduces effective throughput below even the rate-distortion floor). This discrepancy highlights the importance of evidence tiers: Tier~1 and Tier~2 measurements of the ``same'' system can disagree by approximately 0.75 bits.

\section{Discussion}

\subsection{The Receiver-Limited Hypothesis}

The results are consistent with the receiver-limited hypothesis: vocabulary growth reallocates redundancy and sequence length without increasing per-event throughput. The cross-architecture consistency (causal plus encoder-decoder) and the language-match stratification strengthen this claim beyond what a single-architecture sweep could provide. The receiver-limited hypothesis predicts exactly the observed pattern: absolute throughput depends on the receiver's capacity and task, but vocabulary is a redundancy parameter in all cases.

\subsection{Why Biological Systems Cluster at 3--5 Bits}

The convergence is a geometric property of Equation~(4): for $M \in [12, 44]$ and $\varepsilon \in [10^{-4}, 0.05]$, $R_M(\varepsilon)$ maps to 3--5 bits. Biological serial decoders operate in this regime because moderate $M$ reflects the cost-benefit optimum under super-linear discrimination cost, and moderate $\varepsilon$ reflects the thermodynamic floor of reliable discrimination (Hopfield~\cite{hopfield1974}: free energy change $\sim 2$--3 kcal/mol per step). This is a consequence of the rate-distortion surface evaluated at biologically plausible parameters---not an empirical surprise.

\subsection{Vocabulary as Dependent Variable}

Given receiver capacity $C_R$ and error rate $\varepsilon$, the optimal vocabulary size is:
\begin{equation}
    V^* = 2^{C_R + \log_2(1/\varepsilon)}
\end{equation}

Both $C_R$ and $\varepsilon$ are independently measurable, making this a prediction rather than a tautology. This formulation recasts vocabulary size as a derived quantity---a dependent variable determined by receiver architecture and noise characteristics.

\subsection{Cross-Substrate Implications}

The cross-substrate kinship implied by this framework---ribosome, human ear, transformer attention head, and Braille reader all solving the same serial decoding problem---suggests a universal geometric structure governing information processing under noise and cost constraints. This kinship is not by ancestry but by constraint. The mathematics does not distinguish whether the receiver is RNA, neurons, silicon, or raised dots. What matters is the structure of the task: serial, noisy, cost-limited. The broader cross-substrate analysis is presented in the companion paper~\cite{whitmer2026tb}.

\subsection{Falsifiable Predictions}

Four falsifiable predictions are offered:
\begin{enumerate}
    \item \textbf{Vocabulary plateau:} Same architecture trained at $V \in \{256, 1\mathrm{K}, 4\mathrm{K}, 16\mathrm{K}, 64\mathrm{K}, 256\mathrm{K}\}$ should show BPB plateau with no ordinal relationship to $V$.
    \item \textbf{Corruption convergence:} Tokenizers optimized under corruption should converge on smaller, more redundant vocabularies at similar BPB.
    \item \textbf{Genetic code bounds:} All natural genetic codes should encode $15 \leq N \leq 25$ amino acids, consistent with the Tlusty topological bound~\cite{tlusty2008}.
    \item \textbf{Evolutionary convergence:} Evolutionary simulations should converge to 16--32 functional outputs (4--5 bits). Preliminary results from three independent implementations yield $K \sim 19$--30.
\end{enumerate}

\subsection{Limitations}

\begin{enumerate}
    \item \textbf{Near-tautology at low $\varepsilon$.} The ribosome prediction (residual $= 0.0004$ under uniform) is near-tautological because $R_M(\varepsilon) \to \log_2(M)$ as $\varepsilon \to 0$. The framework's value lies in the non-trivial cases.
    \item \textbf{Equiprobable symbol assumption.} All predictions assume symmetric error channels with uniform input. Real systems violate both assumptions. The robustness analysis (Section~3.2) suggests the framework survives this violation but systematic cross-substrate analysis is needed.
    \item \textbf{Observational design.} The 1,749-model experiment is observational, not controlled. Architecture, training data, and training duration confound vocabulary-size effects. A controlled single-architecture sweep is the definitive next step.
    \item \textbf{Tier mixing.} Phonology and music predictions compare $R_M(\varepsilon)$ against support-size bounds [SS], not mutual information [MI]. Conditional entropy rates from Conklin (1995), Temperley (2014), and Pimentel (2020) show music operates at 1.87--2.97 bits/note and phonemes at approximately 3.0 bits---significantly lower than support-size values.
    \item \textbf{AI $I_{\mathrm{eff}}$ is model-dependent}, ranging from approximately 1.4 (frontier) to 5+ (small models). There is no single substrate-level ``AI throughput.''
\end{enumerate}

\section{Conclusion}

The rate-distortion floor $R_M(\varepsilon)$ provides a universal mechanistic bound for serial decoding throughput. For biological receivers with moderate $M$ and moderate $\varepsilon$, this bound confines throughput to 3--5 bits---a geometric property of the rate-distortion surface, not an empirical coincidence. The empirical finding that AI bits-per-byte is independent of vocabulary size across 1,749 models (causal: $p = 0.643$; matched-encoder translation: Spearman $p = 0.13$) supports the vocabulary-as-redundancy hypothesis. Three independent evidence lines---the rate-distortion framework, the biological predictions, and the AI regularities (vocabulary independence and language-match penalty)---together establish a receiver-limited serial decoding constraint. Future work should prioritize controlled single-architecture vocabulary sweeps and Tier~1 confusion-matrix measurements across additional biological systems to narrow the empirical calibration of this framework.

\vspace{1em}
\noindent\textbf{Author Contributions:} Conceptualization, G.L.W.; methodology, G.L.W.; software, G.L.W.; validation, G.L.W.; formal analysis, G.L.W.; investigation, G.L.W.; resources, G.L.W.; data curation, G.L.W.; writing---original draft preparation, G.L.W.; writing---review and editing, G.L.W.; visualization, G.L.W.; supervision, G.L.W.; project administration, G.L.W.

\noindent\textbf{Funding:} This research received no external funding. All work was self-funded by the author.

\noindent\textbf{Data Availability Statement:} All data, analysis scripts, and experiment protocols supporting the findings of this study are publicly available at \url{https://github.com/Windstorm-Labs/receiver-limited-floor} (accessed 12 April 2026).

\noindent\textbf{Acknowledgments:} This paper was improved through adversarial review by six frontier AI models (Claude, GPT, Grok, Gemini, Sonar, and Perplexity). Mathematical derivations and data analysis were performed with the assistance of Claude (Anthropic), an AI research tool. The tokenizer-sweep experiment (1,749 models across two architectures) was executed by Hermes OC1 on an NVIDIA RTX 5090.

\noindent\textbf{Conflicts of Interest:} The author declares no conflict of interest.

\begin{thebibliography}{15}

\bibitem{whitmer2026fons}
Whitmer III, G.L.
The Fons Constraint: Information-Theoretic Convergence on Encoding Depth in Self-Replicating Systems.
\textit{Zenodo} \textbf{2026}.
\href{https://doi.org/10.5281/zenodo.19274048}{DOI: 10.5281/zenodo.19274048}

\bibitem{shannon1948}
Shannon, C.E.
A Mathematical Theory of Communication.
\textit{Bell Syst. Tech. J.} \textbf{1948}, \textit{27}, 379--423.
\href{https://doi.org/10.1002/j.1538-7305.1948.tb01338.x}{DOI: 10.1002/j.1538-7305.1948.tb01338.x}

\bibitem{zaher2009}
Zaher, H.S.; Green, R.
Quality Control by the Ribosome Following Peptide Bond Formation.
\textit{Nature} \textbf{2009}, \textit{457}, 161--166.
\href{https://doi.org/10.1038/nature07582}{DOI: 10.1038/nature07582}

\bibitem{millernicely1955}
Miller, G.A.; Nicely, P.E.
An Analysis of Perceptual Confusions among Some English Consonants.
\textit{J. Acoust. Soc. Am.} \textbf{1955}, \textit{27}, 338--352.
\href{https://doi.org/10.1121/1.1907526}{DOI: 10.1121/1.1907526}

\bibitem{berger1971}
Berger, T.
\textit{Rate Distortion Theory};
Prentice-Hall: Englewood Cliffs, NJ, USA, 1971.

\bibitem{whitmer2026rlf_code}
Whitmer III, G.L.
Receiver-Limited Floor: Experiment Code and Data.
\textit{GitHub} \textbf{2026}.
Available online: \url{https://github.com/Windstorm-Labs/receiver-limited-floor} (accessed 12 April 2026).

\bibitem{tlusty2008}
Tlusty, T.
Rate-Distortion Scenario for the Emergence and Evolution of Noisy Molecular Codes.
\textit{Phys. Rev. Lett.} \textbf{2008}, \textit{100}, 048101.
\href{https://doi.org/10.1103/PhysRevLett.100.048101}{DOI: 10.1103/PhysRevLett.100.048101}

\bibitem{hopfield1974}
Hopfield, J.J.
Kinetic Proofreading: A New Mechanism for Reducing Errors in Biosynthetic Processes Requiring High Specificity.
\textit{Proc. Natl. Acad. Sci. USA} \textbf{1974}, \textit{71}, 4135--4139.
\href{https://doi.org/10.1073/pnas.71.10.4135}{DOI: 10.1073/pnas.71.10.4135}

\bibitem{whitmer2026tb}
Whitmer III, G.L.
The Throughput Basin: Cross-Substrate Convergence and Decomposition of Serial Decoding Throughput.
\textit{Zenodo} \textbf{2026}.
\href{https://doi.org/10.5281/zenodo.19323194}{DOI: 10.5281/zenodo.19323194}

\bibitem{borst1999}
Borst, A.; Theunissen, F.E.
Information Theory and Neural Coding.
\textit{Nat. Neurosci.} \textbf{1999}, \textit{2}, 947--957.
\href{https://doi.org/10.1038/14731}{DOI: 10.1038/14731}

\bibitem{whitmer2026sdb}
Whitmer III, G.L.
The Serial Decoding Basin: Five Experiments on Convergence, Thermodynamic Anchoring, and Receiver-Limited Geometry.
\textit{Zenodo} \textbf{2026}.
\href{https://doi.org/10.5281/zenodo.19323423}{DOI: 10.5281/zenodo.19323423}

\bibitem{whitmer2026dd}
Whitmer III, G.L.
The Dissipative Decoder: Thermodynamic Cost Bounds on the Serial Decoding Throughput Basin.
\textit{Zenodo} \textbf{2026}.
\href{https://doi.org/10.5281/zenodo.19433048}{DOI: 10.5281/zenodo.19433048}

\bibitem{whitmer2026ic}
Whitmer III, G.L.
The Inherited Constraint: Biological Throughput Limits Shape the Information Structure of Human Language and AI.
\textit{Zenodo} \textbf{2026}.
\href{https://doi.org/10.5281/zenodo.19432911}{DOI: 10.5281/zenodo.19432911}

\bibitem{whitmer2026tbo}
Whitmer III, G.L.
The Throughput Basin Origin: Four Orthogonal Experiments on Whether Serial Decoding Convergence Is Architectural, Thermodynamic, or Data-Driven.
\textit{Zenodo} \textbf{2026}.
\href{https://doi.org/10.5281/zenodo.19498582}{DOI: 10.5281/zenodo.19498582}

\bibitem{whitmer2026rlf}
Whitmer III, G.L.
The Receiver-Limited Floor: Rate-Distortion Bounds on Serial Decoding Throughput.
\textit{Zenodo} \textbf{2026}.
\href{https://doi.org/10.5281/zenodo.19322973}{DOI: 10.5281/zenodo.19322973}

\end{thebibliography}

\vspace{2em}
\noindent\textit{This paper is Paper 2 of The Windstorm Series (Papers 1--7) on universal constraints in serial information processing.}

\vspace{1em}
\noindent\textcopyright~2026 Grant Lavell Whitmer III. This work is licensed under a \href{https://creativecommons.org/licenses/by/4.0/}{Creative Commons Attribution 4.0 International License} (CC BY 4.0).

\end{document}
